% TOP_PROBLEM: {"id": 1266, "title": "Lemoine's Conjecture", "category_id": 1, "category": {"id": 1, "name": "number_theory", "display_name": "Number Theory", "description": "Properties of integers, prime numbers, Diophantine equations.", "slug": "number-theory", "order_index": 1, "created_at": "2026-07-31T15:26:25.670Z"}, "status": "open"}
% ATTEMPT_STATUS: unresolved
% Research attempt by GPT_6_Astra_Ultra, 1 October 2026.
% Canonical UnsolvedMath ID; original statements and sources preserved.
% FIRST_POSTED: 2026-10-01
% Imported from: unsolvedmath_additional_500.tex; UnsolvedMath ID 1266
% !TeX program = lualatex
% Compile twice: lualatex 1266.tex
% Self-contained: no external bibliography, images, or \input files.
% Numeric section labels and visible numbers are original UnsolvedMath IDs.
\documentclass[11pt,a4paper]{article}
\usepackage[margin=24mm,headheight=14pt]{geometry}
\usepackage{fontspec}
\setmainfont{Latin Modern Roman}
\setsansfont{Latin Modern Sans}
\setmonofont{Latin Modern Mono}
\usepackage{amsmath,amssymb,mathtools}
\usepackage{unicode-math}
\setmathfont{Latin Modern Math}
\usepackage{microtype}
\usepackage{enumitem,xurl,needspace}
\usepackage[dvipsnames]{xcolor}
\usepackage{hyperref}
\hypersetup{unicode=true,colorlinks=true,linkcolor=MidnightBlue,urlcolor=MidnightBlue,
 pdftitle={UnsolvedMath Problem 1266},pdfauthor={GPT 6 Astra Ultra},bookmarksnumbered=true}
\usepackage{bookmark,tocloft,titlesec,fancyhdr}
\setlength{\cftsecnumwidth}{4.2em}
\cftsetpnumwidth{2.5em}
\cftsetrmarg{4em}
\titleformat{\section}{\Large\bfseries\raggedright}{\thesection}{0.7em}{}
\pagestyle{fancy}\fancyhf{}
\fancyhead[L]{\small\sffamily GPT 6 Astra Ultra research attempt}
\fancyhead[R]{\small Original catalogue IDs}
\fancyfoot[C]{\thepage}
\setlength{\parindent}{0pt}
\setlength{\parskip}{5pt plus 1pt}
\setlength{\emergencystretch}{3em}
\setcounter{tocdepth}{1}\setcounter{secnumdepth}{1}
\setlist[itemize]{leftmargin=3em,itemsep=4pt,topsep=4pt,parsep=0pt}
\allowdisplaybreaks
\newcommand{\N}{\mathbb{N}}
\newcommand{\Z}{\mathbb{Z}}
\newcommand{\Q}{\mathbb{Q}}
\newcommand{\R}{\mathbb{R}}
\newcommand{\C}{\mathbb{C}}
\newcommand{\F}{\mathbb{F}}
\newcommand{\A}{\mathbb{A}}
\newcommand{\PP}{\mathbb{P}}
\newcommand{\E}{\mathbb{E}}
\newcommand{\eps}{\varepsilon}
\newcommand{\dd}{\,\mathrm d}
\newcommand{\sref}[2]{\hyperlink{source:#1:#2}{[S#2]}}
\newif\ifshowcatalogue
\showcataloguetrue
\newcommand{\cataloguescope}[1]{\ifshowcatalogue
 \paragraph{Statement recorded in the supplied source.}
 #1\fi}
\begin{document}
% UnsolvedMath ID 1266; source catalogue code NT-059

\setcounter{section}{1265}

\section{Lemoine's Conjecture}\label{1266}

{\small\sffamily UnsolvedMath ID 1266 · Number Theory}

\subsection{Definitions and mathematical statement}

\paragraph{Shared notation.}Unless a section says otherwise, $\N=\{1,2,\ldots\}$,
$\Z$ is the ring of integers, $\Q,\R,\C$ are the rational, real, and complex
fields, $[N]=\{1,\ldots,\lfloor N\rfloor\}$, and $\log$ is the natural
logarithm. For real functions of a parameter tending to infinity,
$f\ll g$ and $f=O(g)$ mean $|f|\leq Cg$ eventually for some fixed $C$;
$f\gg g$ reverses this comparison; $f=o(g)$ means $f/g\to0$;
$f\sim g$ means $f/g\to1$; and $f\asymp g$ means both order bounds.
A subscript on $O$ or $\ll$ specifies allowed constant dependence.
The expression $n^{o(1)}$ in an upper bound means growth smaller than
$n^\epsilon$ for each fixed $\epsilon>0$. Local definitions govern any
exception, finite-field convention, or use of the same symbol for another
object. An asymptotic claim is not automatically an effective algorithm.



An even semiprime is a product of exactly two primes counting multiplicity, one of which is 2; it therefore has the form $2q$ with $q$ prime. The equation is $n=p+2q$ for primes $p,q$, where $n>5$ is odd. Consequently $p$ is odd; the prime $q$ is allowed to equal 2. \sref{1266}{1} \sref{1266}{2}

\paragraph{Statement recorded in the supplied source.}
Can every odd integer greater than 5 be expressed as the sum of an odd prime and an even semiprime? \sref{1266}{1}

\paragraph{Residual task reported by the archive.}

The embedded review (2026-08-17) records: Prove that every odd n>5 has n=p+2q with p,q prime, or find a counterexample. \sref{1266}{1}

\subsection{Short English statement}

Can every odd integer greater than five be obtained by adding an odd prime to twice a prime?

\subsection{Research attempt}

\paragraph{Scope and route.}
The statement includes $n=7$ and allows $q=2$, so even semiprime four
must not be omitted. Juh\'asz's primary computational paper [S2] was
inspected on 1 October 2026. Its finite-verification setting motivates
an exact counting function, but does not replace control of all odd
integers. The attempt separates local admissibility from the analytic
estimate required for a positive representation count.

\paragraph{An exact finite representation count.}
For an odd integer $N>5$, define
\[
 r(N)=\sum_{\substack{q\le(N-3)/2\\q\ {
m prime}}}
                 1_{N-2q\ {
m prime}}.
\]
The term $p=N-2q$ is then an odd prime at least three. This counts
ordered roles $(p,q)$ rather than unordered pairs, since the coefficient
two distinguishes their roles. The target is exactly $r(N)>0$.
For $N=7$, the pair $(p,q)=(3,2)$ gives the unique representation;
omitting $q=2$ would incorrectly manufacture a counterexample.

For a fixed cutoff $X\ge N$, put
$S_X(\alpha)=\sum_{p\le X}(\log p)e^{2\pi i p\alpha}$.
Orthogonality of the exponentials gives the exact weighted identity
\[
 R(N)=\int_0^1S_X(\alpha)S_X(2\alpha)e^{-2\pi iN\alpha}\,d\alpha
     =\sum_{p+2q=N}\log p\log q.
\]
All weights are positive, so $R(N)>0$ is equivalent to $r(N)>0$.
Replacing either factor by its absolute square changes the equation
being counted and cannot prove positivity of this coefficient.

\paragraph{Local congruences do not forbid a representation.}
For an odd prime $\ell$, choose a nonzero residue $q$ modulo $\ell$
and require $p=N-2q$ also nonzero. If $\ell\mid N$, all $\ell-1$
nonzero choices work. Otherwise exactly one of them, $q=N/2$, is
excluded, leaving $\ell-2$ choices. Both counts are positive. At two,
for odd $N$ the odd-prime model allows $p=q=1$ modulo two, so it too
is admissible. The exceptional small prime choice $q=2$ is handled
by the exact count rather than the all-units local model.

Relative to the independent unit densities these counts produce the
formal local product
\[
 \mathfrak S(N)=2\prod_{\ell>2}\left(1-\frac1{(\ell-1)^2}\right)
               \prod_{\substack{\ell\mid N\\\ell>2}}
                    \frac{\ell-1}{\ell-2}>0.
\]
The infinite product is positive because the sum of its deficits
converges. The archimedean length for $p+2q=N$ is $N/2$, suggesting
a weighted main term $\mathfrak S(N)N/2$. This is a local prediction,
not a derived asymptotic for $R(N)$.

\paragraph{A precise sufficient estimate and the missing cancellation.}
Partition $[0,1]$ into a chosen major-arc set $\mathfrak M$ and its
complement $\mathfrak m$. If one could prove, uniformly for all large
odd $N$, that the major-arc integral equals
$\mathfrak S(N)N/2+o(N)$ and the absolute value of the minor-arc
integral is $o(N)$, positivity would follow from the uniform positive
lower bound on $\mathfrak S(N)$. An explicit threshold plus a finite
check below it would settle the exact target. Neither estimate is
proved in this notebook.

The elementary absolute-value bound does not suffice. Cauchy--Schwarz
and orthogonality give
\[
 \int_0^1|S_X(\alpha)S_X(2\alpha)|\,d\alpha
 \le\sum_{p\le X}(\log p)^2\le X\log^2X.
\]
The equality of the two quadratic norms follows by the change of
variable $\beta=2\alpha$ and periodicity. The resulting bound is
larger than the predicted order-$N$ main term when $X\asymp N$.
It loses the oscillatory cancellation depending on the specified $N$.
An average estimate allowing a sparse exceptional set would also not
prove positivity for every odd integer.

\paragraph{Executed finite verification.}
An integer sieve is used to evaluate $r(N)$ for every odd
$7\le N\le9999$; all counts are positive. Selected exact values are
\[
\begin{array}{c|rrrrrr}
 N&7&9&11&99&999&9999\\\hline
 r(N)&1&2&2&9&39&235
\end{array}
\]
The script checks primality of the actual summands, not merely
congruence admissibility. Files: \texttt{checks\_second.py} and
\texttt{results\_second.json} under
\path{_Local/Erdos_problems/UnsolvedMath/attempt_audit_2026_10_01/number_theory/}.

\paragraph{Remaining gap and outcome.}
The Fourier identity and local counts are exact; their combination
does not by itself yield a positive Fourier coefficient. A pointwise
analytic estimate or another uniform prime-producing argument is
missing. \textbf{Outcome: UNRESOLVED.} The finite range has been
verified, with no assertion about an unbounded tail.

\subsection{Sources}

{\small

\begin{itemize}

\item[\hypertarget{source:1266:1}{S1}] ULAM.AI, \emph{UnsolvedMath}, supplied \texttt{problems.json}, record 1266 (NT-059), statement and background. The embedded literature-review date is 2026-08-17; that is the archive’s review date, not a fresh certification. Dataset landing page: \url{https://huggingface.co/datasets/ulamai/UnsolvedMath}.

\item[\hypertarget{source:1266:2}{S2}] E. Juhász, Empirical Verification of a Generalization of Goldbach's Conjecture, Journal of Integer Sequences 27 (2024). \url{https://cs.uwaterloo.ca/journals/JIS/VOL27/Juhasz/juhasz3.html}. Bibliographic information imported from the supplied record.


\item[E1] E. Juh\'asz, \emph{Empirical Verification of a Generalization
of Goldbach's Conjecture}, Journal of Integer Sequences 27 (2024).
\url{https://cs.uwaterloo.ca/journals/JIS/VOL27/Juhasz/juhasz3.html}.
Primary publication page inspected 1 October 2026. The local check above
is independent and much smaller than published verification work.

\end{itemize}

}

\label{end:1266}
\end{document}
